\par\noindent\textbf{(1)} 设 \(\psi(x,t)=f(x)e^{i\lambda t}\)，\(f\) 为实正且无穷远趋零。代入得 \(\tfrac12f''-\lambda f+f^3=0\)。乘 \(f'\) 并用衰减边界积分得 \((f')^2=2\lambda f^2-f^4\)，要求 \(\lambda>0\)，峰值 \(f_{\max}=\sqrt{2\lambda}\)。
\par\noindent\textbf{(2)} 取 \(f(x)=\eta\,\operatorname{sech}(\eta x)\) 并代入，可得 \(\lambda=\eta^2/2\)。因此 \(\psi=\eta\,\operatorname{sech}(\eta x)e^{i\eta^2t/2}\)，峰值 \(\eta\)，宽度尺度 \(1/\eta\)，说明更强的非线性束缚对应更窄孤子。
\par\noindent\textbf{(3)} 粒子数 \(N=\int|\psi|^2dx=\eta^2\int\operatorname{sech}^2(\eta x)dx=2\eta\)。Hamilton量 \(H=\int[\tfrac12|\psi_x|^2-\tfrac12|\psi|^4]dx\)，利用 \(\int\operatorname{sech}^2u\tanh^2u\,du=2/3\)、\(\int\operatorname{sech}^4u\,du=4/3\)，得到 \(H=-\eta^3/3\)。
\par\noindent\textbf{(4)} Galilean变换 \(\psi_v(x,t)=e^{i(vx-v^2t/2)}\psi(x-vt,t)\) 仍解同一NLS，故亮孤子可匀速传播。若改成散焦号 \(i\psi_t+\psi_{xx}/2-|\psi|^2\psi=0\)，在零背景下无此类正局域亮孤子；应考虑非零背景暗孤子。
